\section{The \sysname{} Framework}
\label{sec:framework}

\sysname{} is a modular experimentation platform for end-to-end,
collocation-aware analysis of ML pipelines. Three principles govern its
design: \emph{modularity and composition}, \emph{declarative configuration},
and \emph{integrated, per-process instrumentation}. It is deliberately
minimalistic: unlike production dataflow engines, it adds no distributed
fault tolerance or autoscaling machinery, because such machinery is exactly
the kind of confound a measurement instrument must not introduce.

\subsection{Core Abstractions}
\label{sec:framework-abstractions}

\paragraph{Stage.} A \emph{stage} is the unit of work and of measurement:
data loading, preprocessing, model inference, retrieval, routing, or custom
logic. A new stage subclasses one base class and implements \texttt{prepare}
(one-time setup: load a model, build an index) and \texttt{run} (per-query
logic). The base class supplies configuration parsing, queue wiring, and
timing instrumentation, so stage authors write only the logic under study.
Each stage executes on its own thread, consuming from input queues and
pushing to output queues; stages with multiple inputs select among
configurable polling policies (wait-for-all merge, first-available,
first-submitted). Queues pass the query object \emph{by reference} --- within
a pipeline, no serialization or copying occurs on the data path
(\S\ref{sec:results-fidelity} quantifies this against a deep-copy counterfactual).

\paragraph{Pipeline.} A \emph{pipeline} is a directed graph of stage
instances --- linear chains, fan-out, fan-in, and \emph{cycles} --- declared
in YAML by listing each stage's identifier and outputs. Cycles express the
retry and rewrite loops of agentic workloads and are bounded by per-query
budgets carried in the query's context. The graph, every stage's parameters,
and the load schedule live in one configuration file; swapping a retriever,
resizing a model, or restructuring the topology is a config edit,
verified by an automatically generated graph visualization.

\paragraph{Load generator.} The \emph{load generator} drives a pipeline with
structured query objects according to a scheduler: a closed-loop scheduler
(one query in flight --- the serial baseline and MLPerf's SingleStream
shape), a Poisson open-loop scheduler (MLPerf Server), a fixed-interval
scheduler (MultiStream), and a saturating all-at-once scheduler (Offline).
Schedulers are components like any other; \S\ref{sec:results-reduction} uses exactly
this pluggability to recover MLPerf's scenario suite. Open-loop schedulers
additionally log every query's \emph{intended} arrival time beside its
realized enqueue time, so coordinated omission --- an entry queue silently
throttling the offered load --- is detected from the trace rather than
assumed away.

\subsection{Execution Model and Isolation}
\label{sec:framework-execution}

A \sysname{} experiment runs in three phases. An orchestrator parses the
configuration and launches \textbf{one OS process per pipeline}; within a
pipeline, stages run as \textbf{threads} sharing one address space. This
split is load-bearing. The process boundary is where isolation and
attribution live: per-process resource listeners (\texttt{macmon} on Apple
silicon; \texttt{nvidia-smi}, \texttt{dcgmi}, \texttt{iostat} on NVIDIA
hosts) sample each pipeline independently, so collocated pipelines ---
declared by simply listing several pipelines in one file --- are separately
attributable. The thread boundary is where zero-copy composition lives:
stages hand off object references, and heavyweight stages release the
interpreter's global lock (GIL) inside native kernels. We are precise about
what overlaps: stages in one process share a default device stream on both
runtimes, so \emph{kernel execution serializes on the accelerator}, and
autoregressive decode re-enters Python every token; cross-stage overlap is
therefore pipelining of CPU-side work (tokenization, retrieval, routing)
against device work and across queries --- scheduling-level concurrency,
not model-level parallelism --- and our metrics and mechanism claims are
worded for exactly that. Pure-Python stages serialize on the GIL; \S\ref{sec:setup-metrics} explains how our concurrency
metric reads under this constraint, and \S\ref{sec:results-fidelity} shows the
hand-off cost is a bounded, depth-independent constant either way.

Robustness follows the same measurement-first logic. A stage that raises
mid-run terminates its pipeline immediately and loudly --- propagating the
stream terminator so the run fails in seconds rather than hanging to a
timeout --- because a silent partial failure is worse than no measurement.
Every run is stamped with the git commit, environment fingerprint, and clock
metadata of the machine that produced it, and the drivers refuse to run on
environments that do not match the pinned specification.

\subsection{Instrumentation}
\label{sec:framework-instrumentation}

Every run emits three artifacts. (i)~A \emph{timing trace}: each stage and
the pipeline driver log start/end events for their prepare and run phases,
stamped with a monotonic clock (\texttt{perf\_counter\_ns}); per-query events
carry the query identifier, submission time, and batch index. Wall-clock
time appears only for cross-process alignment, never in a latency
computation. (ii)~A per-query \emph{answer record} (question, retrieved
documents, generated answer) written by a terminal capture stage, from which
all quality metrics are computed offline. (iii)~The \emph{arrival log}
described above. In parallel, span-level traces link a query's per-stage
slices across threads and processes for timeline visualization, and the
per-process listeners sample power, utilization, and memory. The tracing
layer is optional and its cost is measured, not assumed
(\S\ref{sec:results-fidelity}): with spans exported asynchronously off the
critical path, full instrumentation adds $+2.1\%$ to a multi-millisecond
training step, and can be disabled entirely for minimal-footprint runs.

\subsection{Extensibility}
Stages, schedulers, and polling policies are all resolved from configuration
by import path; adding one means subclassing a base class and naming it in
YAML. The component library used in this paper spans LLM inference across
four backends (HuggingFace, Apple MLX, and served engines), vision
classification and encoding (TorchVision, CLIP, Core~ML/ANE), retrieval
(ChromaDB, FAISS), corpus indexing, the routing and grading stages of
Self-RAG, and synthetic stages (payload injectors, deep-copy counterfactuals,
a calibrated CPU memory-streaming co-runner) used by the fidelity and
contention experiments.
